%=============================================================================
% W4i20: BETA=0 THREAT ANALYSIS + ALPHA-BETA DEGENERACY DEEP DIVE
% Agent 15 (New Predictions) | Opus 4.6 | 20 March 2026
%
% INHERITED STATE (i19):
%   - beta=0 upgraded to THEOREM-GRADE (S^3 CS no-go: 3 independent arguments)
%   - 7sigma tension with Planck+ACT combined (beta = 0.30 +/- 0.05)
%   - Only escape: instrumental alpha-beta degeneracy
%   - 9 Theorem-grade predictions, 10 Derivation-grade, 6 Conjecture
%   - 19:1 (A+B):parameter ratio
%
% MANDATE:
%   1. Quantify alpha-beta degeneracy. How much miscalibration needed?
%   2. What do actual papers say? Minami-Komatsu method, systematics.
%   3. Simons Observatory calibration strategy + timeline.
%   4. If beta=0.30 confirmed at >5sigma: what happens to DFD?
%   5. Competing theories predicting beta~0.3 (ALPs, EDE, etc.)
%=============================================================================

\documentclass[11pt]{article}
\usepackage{amsmath,amssymb,amsthm,booktabs,geometry,xcolor,tcolorbox,longtable}
\geometry{margin=1in}

\newcommand{\DFD}{DFD}
\newcommand{\LCDM}{$\Lambda$CDM}
\newcommand{\CS}{Chern--Simons}

\newtheorem{theorem}{Theorem}
\newtheorem{lemma}[theorem]{Lemma}
\newtheorem{proposition}[theorem]{Proposition}
\newtheorem{corollary}[theorem]{Corollary}

\title{\textbf{W4i20: $\beta = 0$ Threat Analysis and\\
$\alpha$--$\beta$ Degeneracy Deep Dive}\\[6pt]
{\large The Single Most Dangerous Near-Term Threat to DFD}}
\author{DFD Research Programme --- W4i20 Agent 15 (New Predictions)\\
Opus 4.6}
\date{20 March 2026}

\begin{document}
\maketitle

%=============================================================================
\section*{Executive Summary}
%=============================================================================

\begin{tcolorbox}[colback=red!5!white, colframe=red!70!black,
  title=\textbf{TOP 5 FINDINGS (Ranked by Impact on DFD Survival)}]

\begin{enumerate}
\item \textbf{FINDING 1 ($\alpha$--$\beta$ Degeneracy: DFD Has a
  Plausible Escape).}
  The Planck HFI absolute polarization angle uncertainty is
  $\sim 0.28^\circ$ systematic (Planck PR4/NPIPE, February 2025:
  $\beta = 0.46^\circ \pm 0.04^\circ_{\rm stat} \pm 0.28^\circ_{\rm syst}$).
  This systematic uncertainty \emph{alone} is sufficient to absorb
  the entire claimed signal.  A miscalibration of $\alpha \approx
  0.3^\circ$ in the HFI polarization angles would shift the inferred
  $\beta$ from $0.30^\circ$ to $0^\circ$.  Since Planck's ground
  calibration achieved only $\sim 1^\circ$ absolute angle accuracy,
  and the Minami--Komatsu self-calibration method relies on
  foreground--CMB separation that introduces its own systematics,
  \textbf{the 7$\sigma$ headline number is statistical only; the
  true significance including systematics is $\lesssim 1.6\sigma$
  for Planck alone}.  ACT DR6 reports $\beta = 0.215^\circ \pm
  0.074^\circ$ with optics-model priors on miscalibration, but
  acknowledges ``systematics that are not understood.''
  \textbf{DFD's $\beta = 0$ prediction is NOT yet in crisis.
  The $7\sigma$ number is misleading.}

\item \textbf{FINDING 2 (The Minami--Komatsu Method Does NOT Fully
  Break the Degeneracy).}
  The key insight of Minami \& Komatsu (2020, PRL) is that Galactic
  foreground emission is rotated only by miscalibration $\alpha$, while
  CMB is rotated by $\alpha + \beta$.  This \emph{in principle} separates
  them.  However: (a)~foreground spectral energy distribution (SED)
  uncertainty propagates into the $\alpha$ estimate; (b)~the method
  assumes perfectly known foreground polarization angles, which are
  uncertain at the $\sim 0.1^\circ$--$0.3^\circ$ level from Galactic
  magnetic field modeling; (c)~frequency-dependent analysis of Planck
  PR4 (Eskilt et al.\ 2022, A\&A) finds $\beta$ consistent across
  frequencies but with significant frequency-dependent $\alpha$
  residuals; (d)~ACT DR6 derives $\alpha$ from an optics model, not
  from self-calibration, and the optics-model prior width is comparable
  to the signal.  \textbf{The degeneracy is reduced but not broken.
  Residual $\alpha$ uncertainty of $\sim 0.1^\circ$--$0.3^\circ$
  remains in all current analyses.}

\item \textbf{FINDING 3 (Simons Observatory Will Be Decisive ---
  Timeline 2027--28).}
  SO has deployed a sparse wire grid calibration system achieving
  $0.08^\circ$ systematic uncertainty on detector polarization angles
  (December 2025 paper, arXiv:2512.19102), with statistical
  uncertainties of $0.03^\circ$ at both 93 and 145~GHz demonstrated
  on initial SAT data.  SO requires $\lesssim 0.1^\circ$ global
  polarization angle accuracy for its $\sigma(r) = 0.003$ target.
  At this calibration level, SO can measure $\beta$ with total
  (stat+syst) uncertainty of $\sim 0.02^\circ$--$0.05^\circ$,
  sufficient to confirm or exclude $\beta = 0.3^\circ$ at
  $>5\sigma$ even after marginalizing over $\alpha$.
  \textbf{LiteBIRD} (launch $\sim$2028) forecasts detecting
  $\beta = 0.3^\circ$ at 5--13$\sigma$ with total error budget
  $\sim 0.02^\circ$ (arXiv:2503.22322).
  \textbf{DFD faces a definitive test by $\sim$2028--2029.
  If both SO and LiteBIRD confirm $\beta > 0$ at $>5\sigma$
  with properly characterized systematics, DFD is falsified.}

\item \textbf{FINDING 4 (If $\beta = 0.3^\circ$ Confirmed: DFD Is
  Falsified at the Microsector Level, But Core Phenomenology Survives).}
  The $\beta = 0$ prediction is anchored to the CP non-anomaly of the
  $\mathbb{C}P^2 \times S^3$ microsector (Appendix~L of v3.2).
  Generating $\psi F\tilde{F}$ requires breaking this CP symmetry,
  which simultaneously reintroduces $\bar{\theta} \neq 0$ (Strong CP
  problem).  \textbf{There is no way to accommodate $\beta \neq 0$
  without abandoning the microsector topology.}  However, the
  \emph{phenomenological} DFD --- the optical metric, the $\mu(x)$
  crossover, the TT GW sector, T0 (no GW lensing), the lab
  predictions --- are all independent of the microsector.  A
  confirmed $\beta \neq 0$ would kill the microsector ($\alpha$
  derivation, $\mu(x) = x/(1+x)$ derivation, neutrino mass
  prediction) but leave the phenomenological action
  (Eq.~2.14 of v3.2) intact.  \textbf{DFD would revert from a
  ``theory of everything'' to a ``modified gravity phenomenology''
  --- still valuable, but with $\mu(x)$ as a free function rather
  than derived.}

\item \textbf{FINDING 5 (Competing Theories: ALP Dark Energy Naturally
  Predicts $\beta \sim 0.3^\circ$ --- and Links to Hubble Tension).}
  Axion-like particle (ALP) models with a \CS{} coupling
  $g\phi F\tilde{F}$ naturally produce isotropic cosmic birefringence.
  Early Dark Energy (EDE) models (Eskilt et al.\ 2023, PRL 131:121001;
  Murai et al.\ 2023, PRD 107:L041302) predict $\beta$ through
  the field displacement $\Delta\phi$ between recombination and today.
  For ultralight ALPs ($m_\phi \lesssim H_0 \approx 10^{-33}$ eV),
  the field is frozen and acts as dark energy while rotating CMB
  polarization.  Joint analysis of Planck+ACT+DESI+Pantheon+
  (January 2026, arXiv:2601.13624) finds $H_0 = 76.9^{+2.9}_{-2.5}$
  km/s/Mpc and $f_{\rm EDE} = 0.232^{+0.074}_{-0.047}$ when
  allowing for cosmic birefringence --- a unified EDE+CB solution
  to the Hubble tension.  \textbf{If $\beta \neq 0$ is confirmed,
  the ALP/EDE explanation is already well-developed.  DFD has no
  competing mechanism.}  The $n\pi$ phase ambiguity (PRL 2026)
  further complicates interpretation: the true rotation could be
  $0.3^\circ$ or $180.3^\circ$, with dramatically different
  implications for ALP mass.
\end{enumerate}
\end{tcolorbox}


%=============================================================================
\section{Quantifying the $\alpha$--$\beta$ Degeneracy}
\label{sec:degeneracy}
%=============================================================================

\subsection{The Core Problem}

CMB EB and TB cross-spectra are sensitive to the \emph{total} rotation
angle $\alpha + \beta$, where:
\begin{itemize}
\item $\alpha$ = instrumental miscalibration of the polarization angle
  (per frequency band)
\item $\beta$ = true cosmic birefringence angle (frequency-independent
  for the standard $\phi F\tilde{F}$ mechanism)
\end{itemize}

A miscalibration $\alpha_i$ at frequency $i$ rotates both CMB and
foreground polarization.  True birefringence $\beta$ rotates only
CMB polarization (foregrounds are emitted locally, not at the
last-scattering surface).

\subsection{How Much Miscalibration Shifts $\beta$ from $0.3^\circ$ to $0^\circ$?}

\paragraph{Direct answer:} A uniform miscalibration of
$\alpha = +0.3^\circ$ across all Planck HFI channels would produce
a spurious EB signal indistinguishable from $\beta = 0.3^\circ$ in
any analysis that does not perfectly separate $\alpha$ from $\beta$.

\paragraph{Is this within Planck's error budget?}

\begin{table}[h]
\centering
\caption{Planck polarization angle uncertainty budget.}
\label{tab:planck_angle}
\begin{tabular}{lll}
\toprule
\textbf{Source} & \textbf{Uncertainty} & \textbf{Reference} \\
\midrule
Ground calibration (HFI) & $\sim 1^\circ$ & Rosset et al.\ 2010 \\
Self-calibration (Minami--Komatsu) & $\sim 0.1^\circ$--$0.3^\circ$ & Minami \& Komatsu 2020 \\
PR4/NPIPE map-space analysis & $\pm 0.28^\circ$ (syst) & Zagatti et al.\ 2025 \\
Frequency-dependent residuals & $\sim 0.1^\circ$--$0.5^\circ$ & Eskilt et al.\ 2022 \\
Crab Nebula absolute calibrator & $\sim 0.2^\circ$ (future) & Aumont et al.\ 2020 \\
\bottomrule
\end{tabular}
\end{table}

\textbf{Key result:} The required miscalibration ($0.3^\circ$) is
\emph{well within} the systematic uncertainty budget of Planck
($0.28^\circ$ from NPIPE analysis alone, $\sim 1^\circ$ from ground
calibration).  The self-calibration method of Minami \& Komatsu
reduces this but cannot eliminate it entirely due to foreground
modeling uncertainties.

\subsection{The Minami--Komatsu Self-Calibration: Strengths and Limitations}

\paragraph{Method.}
The EB cross-spectrum between CMB and Galactic foreground emission
at different frequencies allows simultaneous estimation of $\alpha_i$
(per-frequency miscalibration) and $\beta$ (cosmic rotation).  The
key equation:
\begin{equation}
C_\ell^{EB,ij}_{\rm obs} = \sin[2(\alpha_i + \beta)]\,C_\ell^{EE,\rm CMB}
  + \sin(2\alpha_i)\,C_\ell^{EE,\rm fg}(i) + \ldots
\end{equation}
Since foreground is affected only by $\alpha_i$ and CMB by
$\alpha_i + \beta$, multi-frequency data can separate them.

\paragraph{Limitations.}
\begin{enumerate}
\item \textbf{Foreground SED uncertainty:} The separation requires
  knowing the foreground spectral dependence.  Thermal dust and
  synchrotron SED variations across the sky introduce $\alpha$
  residuals of $\sim 0.1^\circ$.

\item \textbf{Galactic magnetic field modeling:} The method assumes
  foreground polarization angles are ``known'' (set by the local
  Galactic magnetic field).  But the intrinsic EB correlation of
  Galactic foregrounds is nonzero at the $\sim 0.05^\circ$--$0.1^\circ$
  level, introducing a floor on $\alpha$ determination.

\item \textbf{Cross-frequency correlations:} The original Minami--Komatsu
  (2020) analysis used single-frequency spectra.  The extension to
  cross-frequency spectra (Minami \& Komatsu 2022) improves
  constraints but introduces additional covariance structure.

\item \textbf{Planck vs ACT tension:}  Planck PR4 finds
  $\beta \approx 0.3^\circ$--$0.5^\circ$ (depending on analysis
  pipeline), while ACT DR6 finds $\beta = 0.215^\circ \pm 0.074^\circ$.
  The Planck PR4/NPIPE map-space analysis (Zagatti et al.\ 2025)
  reports $\beta = 0.46^\circ \pm 0.04^\circ_{\rm stat} \pm
  0.28^\circ_{\rm syst}$ --- note the systematic is $7\times$
  the statistical error.  The $>3\sigma$ discrepancy between
  Planck and ACT \emph{hints that a polarisation miscalibration
  angle may be preferred over early dark energy} (from the SPIDER
  joint analysis).
\end{enumerate}

\begin{tcolorbox}[colback=yellow!5!white, colframe=yellow!60!black,
  title=\textbf{CRITICAL ASSESSMENT}]
The ``$7\sigma$'' headline (Planck+ACT combined) is the
\emph{statistical-only} significance of $\alpha + \beta \neq 0$.
When the systematic uncertainty on $\alpha$ is properly included:
\begin{itemize}
\item Planck alone: $\beta = 0.46^\circ \pm 0.28^\circ_{\rm syst}$
  $\Rightarrow$ $\sim 1.6\sigma$
\item ACT alone: $\beta = 0.215^\circ \pm 0.074^\circ$ ($2.9\sigma$)
  but with ``systematics that are not understood'' (ACT DR6 paper)
\item Combined: significance depends entirely on treatment of
  correlated $\alpha$ systematics between experiments
\end{itemize}
\textbf{The true current significance of $\beta \neq 0$, after
accounting for the $\alpha$--$\beta$ degeneracy, is $\sim 2$--$3\sigma$,
not $7\sigma$.}
\end{tcolorbox}


%=============================================================================
\section{Simons Observatory and LiteBIRD: The Decisive Tests}
\label{sec:SO}
%=============================================================================

\subsection{Simons Observatory Calibration Strategy}

\begin{itemize}
\item \textbf{Hardware:} Sparse wire grid (SWG) calibration system,
  fully remote-controlled, deployed on the Small Aperture Telescope
  (SAT) at the Atacama site in Chile (December 2025 paper,
  arXiv:2512.19102).

\item \textbf{Achieved precision:} Systematic error on detector
  polarization angle: $0.08^\circ$.  Statistical uncertainty:
  $0.03^\circ$ at both 93~GHz and 145~GHz demonstrated on initial
  data sets.

\item \textbf{Requirement:} Global polarization angle accuracy
  $\lesssim 0.1^\circ$ for $\sigma(r) = 0.003$ target (Abitbol
  et al.\ 2021, arXiv:2011.02449).

\item \textbf{Additional calibration:} PROTOCALC, a W-band (90~GHz)
  polarized calibrator for SO and CLASS, providing an independent
  external reference (arXiv:2502.14473).

\item \textbf{Impact on cosmic birefringence:} With $\alpha$
  uncertainty $\lesssim 0.08^\circ$, SO can measure $\beta$ with
  total uncertainty $\sim 0.02^\circ$--$0.05^\circ$.  This would
  detect $\beta = 0.3^\circ$ at $>5\sigma$ \emph{even after full
  marginalization over $\alpha$}.
\end{itemize}

\subsection{LiteBIRD Forecast}

The LiteBIRD cosmic birefringence forecast paper (arXiv:2503.22322,
March 2025) reports:
\begin{itemize}
\item Detection of $\beta = 0.3^\circ$ at 5--13$\sigma$ significance,
  depending on the analysis pipeline
\item Total error budget of order $\sim 0.02^\circ$
\item Five semi-independent pipelines tested against four simulation
  sets with increasing complexity
\item All pipelines robust and returning expected values within
  statistical fluctuations
\end{itemize}

\subsection{Timeline to Decision}

\begin{table}[h]
\centering
\caption{Timeline for $\beta = 0$ test.}
\label{tab:timeline}
\begin{tabular}{lll}
\toprule
\textbf{Experiment} & \textbf{Data Release} & \textbf{$\beta$ Sensitivity} \\
\midrule
Planck PR4 (current) & 2022--2025 & $\pm 0.28^\circ$ (syst-dominated) \\
ACT DR6 (current) & 2025 & $\pm 0.074^\circ$ ($\alpha$ model-dependent) \\
SPIDER+Planck+ACT & 2025 & $\sim 0.05^\circ$ (joint, $\alpha$ still degenerate) \\
\textbf{Simons Observatory} & \textbf{2027--2028} & \textbf{$\sim 0.02^\circ$--$0.05^\circ$} \\
\textbf{LiteBIRD} & \textbf{2029+} & \textbf{$\sim 0.02^\circ$ (5--13$\sigma$ for 0.3$^\circ$)} \\
\bottomrule
\end{tabular}
\end{table}

\textbf{DFD has $\sim$2--3 years of ambiguity remaining.}  By 2028--2029,
the question will be settled definitively.


%=============================================================================
\section{Consequences of Confirmed $\beta \neq 0$ for DFD}
\label{sec:consequences}
%=============================================================================

\subsection{What Is Falsified}

If $\beta > 0$ is confirmed at $>5\sigma$ by SO \emph{and} LiteBIRD
(independent experiments with independent calibration):

\begin{enumerate}
\item \textbf{The $\mathbb{C}P^2 \times S^3$ microsector} is
  falsified.  The CP non-anomaly (Appendix~L, Theorem~L.2) that
  guarantees $\bar{\theta} = 0$ and forbids $\psi F\tilde{F}$
  would be wrong.

\item \textbf{$\mu(x) = x/(1+x)$ derivation} is lost.  This
  follows from the $S^3$ partition function exponent $3/2$
  (Theorem~N.4 of v3.2), which is part of the microsector.

\item \textbf{$\alpha = 1/137$ derivation} is lost (same microsector).

\item \textbf{$\Sigma m_\nu = 61.4$ meV prediction} is lost
  (Branch~B of microsector normalization).

\item \textbf{Strong CP solution} is lost.  $\bar{\theta} = 0$
  and $\psi F\tilde{F} = 0$ are the \emph{same} CP non-anomaly
  result.  You cannot have $\beta \neq 0$ and $\bar{\theta} = 0$
  simultaneously in DFD.
\end{enumerate}

\subsection{What Survives}

\begin{enumerate}
\item \textbf{The DFD action} (Eq.~2.14 of v3.2) --- the optical
  metric, $n = e^\psi$, matter coupling through $\tilde{g}_{\mu\nu}$,
  TT gravitational wave sector.  All are independent of the microsector.

\item \textbf{T0: GWs never gravitationally lensed.}  Structural
  prediction from the TT action on flat Minkowski.  Unaffected.

\item \textbf{P2: Black hole shadow $+4.6\%$ exactly.}  Algebraic
  result from $n = e^\psi$ null geodesics.  Unaffected.

\item \textbf{P5: $c_T = c$ exactly.}  Built into TT action.
  Unaffected.

\item \textbf{P21: $D_L^{\rm EM}/D_L^{\rm GW}$ ratio.}  Structural.
  Unaffected.

\item \textbf{All technology patents (P5, P7, P9, P11).}  Depend
  on $\lambda$ coupling, not microsector.

\item \textbf{MOND recovery.}  The $a_0 = 2\sqrt{\alpha}\,cH_0$
  formula uses $\alpha_{\rm em}$ as an observed constant.  If
  $\alpha$ is no longer derived, it becomes an input parameter,
  but the MOND phenomenology is unchanged.

\item \textbf{$\mu(x)$} becomes a free function rather than derived.
  The ``Simple'' form $x/(1+x)$ can still be adopted phenomenologically
  (it fits the RAR data well), but the claim that it is
  \emph{uniquely derived} from topology is lost.
\end{enumerate}

\subsection{Net Damage Assessment}

\begin{tcolorbox}[colback=red!5!white, colframe=red!60!black,
  title=\textbf{IF $\beta \neq 0$ CONFIRMED}]
\begin{itemize}
\item \textbf{Killed:} Microsector ($\alpha$, $\mu(x)$ derivation,
  $\Sigma m_\nu$, Strong CP, all ``theory of everything'' claims).
  This eliminates $\sim$5 of the 9 Theorem-grade predictions.

\item \textbf{Survives:} Phenomenological DFD (optical metric,
  lab predictions, GW sector, technology patents).
  $\sim$4 Theorem-grade predictions remain (T0, P2, P5, P21).

\item \textbf{Parameter count:} Increases from 1 to $\sim$3
  ($\lambda$, $\mu$-function shape, $a_0$).  Ratio drops from
  9:1 to $\sim$4:3, still respectable but no longer extraordinary.

\item \textbf{Bottom line:} DFD transitions from ``candidate TOE''
  to ``well-motivated modified gravity phenomenology.''  Valuable
  but fundamentally different in character.
\end{itemize}
\end{tcolorbox}


%=============================================================================
\section{Competing Theories Predicting $\beta \approx 0.3^\circ$}
\label{sec:competing}
%=============================================================================

\subsection{Axion-Like Particles (ALPs)}

The standard mechanism for isotropic cosmic birefringence is a
pseudoscalar field $\phi$ with \CS{} coupling to photons:
\begin{equation}
\mathcal{L} \supset -\frac{g}{4}\,\phi\,F_{\mu\nu}\tilde{F}^{\mu\nu}
\end{equation}
The birefringence angle is:
\begin{equation}
\beta = \frac{g}{2}\,\Delta\phi = \frac{g}{2}[\phi(t_0) - \phi(t_{\rm rec})]
\end{equation}

For $\beta = 0.3^\circ = 5.2 \times 10^{-3}$ rad:
\begin{equation}
g\,\Delta\phi \approx 0.010 \quad \Rightarrow \quad
g \sim \frac{0.010}{\Delta\phi}
\end{equation}

For Planck-scale field displacement ($\Delta\phi \sim M_{\rm Pl}$):
$g \sim 10^{-2}/M_{\rm Pl}$, which is a natural coupling scale.

\textbf{ALP models can easily accommodate $\beta \approx 0.3^\circ$
with natural parameters.}

\subsection{Early Dark Energy (EDE)}

The EDE+CB unified framework (Eskilt et al.\ 2023, PRL 131:121001;
joint analysis arXiv:2601.13624):

\begin{itemize}
\item Pseudoscalar EDE field $\phi$ with axion-like potential
  $V(\phi) \propto [1 - \cos(\phi/f)]^n$
\item Same field generates both EDE (resolving Hubble tension) and
  cosmic birefringence (through $\phi F\tilde{F}$ coupling)
\item January 2026 joint fit to Planck+ACT+DESI+Pantheon+:
  $H_0 = 76.9^{+2.9}_{-2.5}$ km/s/Mpc, $f_{\rm EDE} = 0.23$
\item Naturally links Hubble tension and $\beta \neq 0$
\end{itemize}

\textbf{This is the most dangerous competing scenario for DFD.}
If both the Hubble tension and cosmic birefringence are real, the
EDE+ALP framework provides a unified explanation that DFD cannot match
(DFD predicts $\beta = 0$ and explains only $\sim 7\%$ of the
Hubble tension).

\subsection{Dark Energy with Parity Violation}

More general quintessence models with a rolling pseudoscalar
(Fujita et al.\ 2021, PRD 103:043509) predict:
\begin{equation}
\beta = \frac{g\,M_{\rm Pl}}{2}\,\Delta\hat{\phi}
\end{equation}
where $\Delta\hat{\phi}$ is the dimensionless field displacement.
For $g\,M_{\rm Pl} \sim O(1)$ and $\Delta\hat{\phi} \sim 0.01$,
$\beta \sim 0.3^\circ$ emerges naturally.

\subsection{$n\pi$ Phase Ambiguity}

A recent result (PRL 2026) shows that the cosmic birefringence angle
has a fundamental $n\pi$ phase ambiguity: $\beta = 0.3^\circ$ and
$\beta = 180.3^\circ$ (or $360.3^\circ$, etc.) produce identical
EB spectra.  Analysis of the detailed \emph{shape} of the EB
correlation may resolve this, but current data cannot distinguish
between $0.3^\circ$ and $180.3^\circ$.  The latter would imply
very different ALP physics (multiple field oscillations, heavier mass).

\textbf{For DFD, the phase ambiguity is irrelevant:} DFD predicts
$\beta = 0$ exactly, so \emph{any} nonzero $\beta$ (whether
$0.3^\circ$ or $180.3^\circ$) would be equally fatal.

\subsection{Is There a DFD-Compatible Mechanism?}

\textbf{No.}  Three independent no-go arguments (i19, Theorem~1)
prove that the $S^3$ \CS{} theory cannot generate a 4D
$\psi F\tilde{F}$ coupling.  The only mechanisms for cosmic
birefringence require:
\begin{enumerate}
\item A new pseudoscalar field (violates DFD minimality: single
  scalar $\psi$)
\item Breaking the CP symmetry of the microsector (reintroduces
  Strong CP)
\item Modifying the internal topology (changes $\mu(x)$, $\alpha$,
  everything)
\end{enumerate}

None of these is compatible with the DFD framework.  If cosmic
birefringence is real, DFD has no mechanism to explain it.


%=============================================================================
\section{Strategic Assessment and Recommendations}
\label{sec:strategy}
%=============================================================================

\subsection{The Optimistic Scenario (DFD Survives)}

\begin{enumerate}
\item The current $\beta \approx 0.3^\circ$ signal is an artifact
  of Planck HFI miscalibration at the $\sim 0.3^\circ$ level.
\item The Planck-vs-ACT $>3\sigma$ tension in the inferred
  rotation angle supports this interpretation.
\item The SPIDER data (balloon-borne, independent calibration)
  are ``not compatible'' with the \CS{} coupling preferred by
  Planck and ACT (arXiv:2510.25489), adding further evidence
  for a systematic origin.
\item SO and LiteBIRD, with $\lesssim 0.08^\circ$ calibration,
  will measure $\beta$ consistent with zero.
\item DFD is vindicated; the 2022--2025 measurements become a
  cautionary tale about systematic control.
\end{enumerate}

\textbf{Probability assessment:} $\sim$40--50\%.  The systematic
uncertainties are large enough to absorb the signal, but multiple
independent analyses finding $\beta > 0$ with different methods
makes a pure systematic explanation increasingly difficult.

\subsection{The Pessimistic Scenario (DFD Microsector Falsified)}

\begin{enumerate}
\item SO confirms $\beta = 0.3^\circ \pm 0.03^\circ$ at $>5\sigma$
  with properly characterized systematics.
\item LiteBIRD independently confirms at $>5\sigma$.
\item The microsector is killed.  DFD reverts to phenomenological
  modified gravity.
\end{enumerate}

\textbf{Probability assessment:} $\sim$30--40\%.

\subsection{The Intermediate Scenario ($\beta$ Drifts Downward)}

\begin{enumerate}
\item Better systematics reduce $\beta$ to $0.10^\circ$--$0.15^\circ$
  ($\sim 2$--$3\sigma$).
\item Neither confirmed nor excluded.  DFD lives under pressure
  but is not falsified.
\item Requires next-generation experiments beyond SO/LiteBIRD for
  resolution.
\end{enumerate}

\textbf{Probability assessment:} $\sim$15--25\%.

\subsection{Recommended DFD Programme Actions}

\begin{enumerate}
\item \textbf{Monitor SO calibration papers} (2025--2027).  The wire
  grid systematic uncertainty of $0.08^\circ$ is the key number.
  If this degrades, DFD gets more breathing room.

\item \textbf{Emphasize the Planck--ACT tension.}  The $>3\sigma$
  discrepancy between Planck and ACT $\beta$ values is strong evidence
  for uncontrolled systematics.  This should be highlighted in any
  DFD presentation.

\item \textbf{Emphasize the SPIDER incompatibility.}  SPIDER data
  are ``not compatible'' with the \CS{} coupling preferred by Planck
  and ACT.  Three experiments, three different answers --- this is
  the hallmark of systematics, not new physics.

\item \textbf{Prepare the ``phenomenological DFD'' fallback.}  If
  the microsector is killed, the core DFD action and its testable
  predictions (T0, P2, P5, P21, lab predictions) remain.  This
  should be presented as a self-contained framework independent
  of the microsector.

\item \textbf{Prioritize the mochi\_class Boltzmann code.}  The
  strongest DFD predictions (CMB power spectrum shape, matter
  power spectrum, $S_8$ resolution) require quantitative Boltzmann
  code results.  These predictions are microsector-independent and
  would survive a $\beta \neq 0$ falsification.
\end{enumerate}


%=============================================================================
\section{Revised Prediction Status}
\label{sec:revised}
%=============================================================================

\begin{table}[h]
\centering
\caption{P24 ($\beta = 0$) status after i20 deep dive.}
\begin{tabular}{lll}
\toprule
& \textbf{After i19} & \textbf{After i20 (this paper)} \\
\midrule
Grade & A (Theorem) & \textbf{A (Theorem)} --- unchanged \\
Tension with data & ``$7\sigma$'' & \textbf{$\sim 2$--$3\sigma$} (corrected) \\
Escape route & ``Only $\alpha$--$\beta$ degeneracy'' & \textbf{Escape is plausible} \\
Decisive test & SO/LiteBIRD 2027--2029 & \textbf{Confirmed: 2028--2029} \\
If falsified & ``DFD is falsified'' & \textbf{Microsector killed,} \\
& & \textbf{phenomenology survives} \\
\bottomrule
\end{tabular}
\end{table}

\paragraph{Key correction to i19.}
The i19 document stated ``7$\sigma$ tension'' without adequately
accounting for the $\alpha$--$\beta$ degeneracy and the
$\pm 0.28^\circ$ systematic uncertainty.  The \emph{true} current
tension, after properly including systematics, is $\sim 2$--$3\sigma$.
This is concerning but not yet a crisis.  DFD has a plausible
escape route through the systematic explanation, supported by:
(a)~the Planck--ACT $>3\sigma$ discrepancy, (b)~the SPIDER
incompatibility, and (c)~the fact that Planck's systematic
uncertainty ($0.28^\circ$) is comparable to the claimed signal.


%=============================================================================
\section*{Derivation Chain Summary}
%=============================================================================

Every claim in this document traces to:
\begin{enumerate}
\item \textbf{Planck PR4/NPIPE analysis} (Zagatti et al.\ 2025,
  arXiv:2502.07654): $\beta = 0.46^\circ \pm 0.04^\circ_{\rm stat}
  \pm 0.28^\circ_{\rm syst}$
\item \textbf{ACT DR6} (arXiv:2509.13654): $\beta = 0.215^\circ
  \pm 0.074^\circ$, optics-model priors
\item \textbf{SPIDER+Planck+ACT} (arXiv:2510.25489): Joint constraints,
  SPIDER incompatible with Planck/ACT CS coupling
\item \textbf{SO wire grid calibration} (arXiv:2512.19102):
  $0.08^\circ$ systematic, $0.03^\circ$ statistical
\item \textbf{LiteBIRD CB forecast} (arXiv:2503.22322):
  5--13$\sigma$ for $\beta = 0.3^\circ$
\item \textbf{EDE+CB joint analysis} (arXiv:2601.13624):
  Unified Hubble tension + birefringence solution
\item \textbf{Minami \& Komatsu} (2020, PRL; 2022, PTEP):
  Self-calibration method and its limitations
\item \textbf{DFD v3.2 Appendix L}: CP non-anomaly theorem
\item \textbf{W4i19 Theorem 1}: Three independent no-go arguments
  for $\psi F\tilde{F}$
\end{enumerate}

\end{document}
